\documentclass[10pt]{article}
\usepackage{tmlr}

\input{math_commands.tex}

\usepackage{hyperref}
\usepackage{url}
\usepackage{booktabs}
\usepackage{multirow}
\usepackage{amsfonts}
\usepackage{amssymb}
\usepackage{graphicx}

\title{The Safety Metric is the Independent Variable:\\
Pairwise Reversal Rates Reveal What Aggregate Scores Hide\\
in LLM Safety Evaluation}

\author{\name Anonymous \email anon@example.com \\
       \addr Anonymous Affiliation}

\def\month{05}
\def\year{2026}
\def\openreview{\url{https://openreview.net/forum?id=ANONYMOUS}}

\begin{document}
\maketitle

\begin{abstract}
Frontier-LLM evaluation has standardised on a ``safety SOTA'' framing
in which a single benchmark --- HarmBench, AlpacaEval, AIR-Bench, BBQ
--- serves as the headline number. Stanford's HEIM-style holistic
evaluations (HELM Safety v1.17.0, 2025) acknowledge that ``no single
model excels in all aspects'' but the case-level cost of metric choice
has not been quantified. We import the criterion-as-independent-
variable framework, validated on LLM Q\&A, fund ranking, and
text-to-image evaluation in companion papers, and apply it to the 87
frontier-LLM~$\times$~9 safety-metric panel released by HELM Safety
v1.17.0. We compute the pairwise reversal-rate matrix with cluster-
bootstrap CIs ($B = 1000$ over models). Headline findings:
\textbf{(1)} four metric pairs exceed 50\,\% reversal --- i.e., the
preferred LLM flips on more than half of model pairs when the metric
is swapped --- with the maximum at SimpleSafetyTests~$\leftrightarrow$~
BBQ-unambiguous reaching \textbf{63.8\,\% (95\,\% CI [54.9\,\%,
71.7\,\%])}; \textbf{(2)} every harm/safety metric exhibits
\emph{negative} Spearman~$\rho$ with BBQ-unambiguous bias accuracy
($\rho \in [-0.34, -0.10]$), so the LLMs that score highest on safety
benchmarks score \emph{worst} on the bias benchmark; \textbf{(3)}
within-HarmBench across-annotator reversal between the LM-evaluated
and Llama-annotator scores is 25\,\%, demonstrating that ``HarmBench
score'' is not a single quantity even within one benchmark;
\textbf{(4)} hierarchical clustering on $1 - \rho$ yields three
clusters that align with the underlying axis structure (harm/refusal,
bias-favouring, bias-disfavouring). We argue the choice of safety
metric is an \textbf{independent variable} of the evaluation research
design, not a measurement nuisance, and recommend that
LLM-developer ``safety SOTA'' claims report the full reversal-rate
matrix alongside their headline.
\end{abstract}

\section{Introduction}

When OpenAI, Anthropic, Google, or Meta announce a new LLM, the
accompanying safety claim is typically a single benchmark number:
``refuses 98\,\% of HarmBench prompts,'' ``scores 0.95 on Anthropic
Red-Teaming,'' ``best-in-class on AIR-Bench.'' These are non-trivially
valuable measurements. They also conceal a structural problem: the
safety landscape has accumulated 102+ benchmarks since 2018, and the
benchmarks \textbf{disagree}.

Stanford's HEIM Safety \citep{liang2024helm} and AIR-Bench
\citep{zeng2024airbench} document the disagreement qualitatively.
BEAVERTAILS \citep{ji2023beavertails} showed that helpfulness and
harmlessness can diverge per-prompt. \citet{liu2025scalesofjustitia}
catalogue over 100 benchmarks and conclude the field is fragmented.
None of these works has reported the \emph{case-level statistic}
that practitioners actually care about: when I switch from one
benchmark to another, how often does my model's ranking flip?

We answer that question. Using the public HELM Safety v1.17.0
release (87 frontier LLMs $\times$ 9 safety metrics across 5
benchmarks: HarmBench, SimpleSafetyTests, XSTest, BBQ,
AnthropicRedTeam), we compute the full pairwise reversal-rate matrix
with cluster-bootstrap CIs and Spearman~$\rho$.

\subsection{Contributions}

\textbf{(C1)} A pre-registered statistical test of metric
disagreement on LLM safety: 87-LLM panel, 9 metrics, 1000-trial
cluster bootstrap.

\textbf{(C2) Negative correlation between safety and bias.} Every
harm/refusal metric on HELM Safety exhibits negative
Spearman~$\rho$ with the BBQ-unambiguous bias accuracy
($\rho \in [-0.345, -0.101]$). The LLMs that score highest on
harm/refusal benchmarks score \emph{worst} on the bias benchmark.

\textbf{(C3) Within-benchmark instability.} Even within one
benchmark (HarmBench), the LM-evaluated safety score and the
Llama-annotator success rate produce different rankings on
25.5\,\% of LLM pairs. The ``HarmBench score'' is not a single
quantity.

\textbf{(C4) Three-cluster axis structure.} Hierarchical clustering
on $1 - \rho$ separates the metrics into a harm/refusal main
cluster, a bias contexts cluster, and a saturated outlier. The axes
are nearly orthogonal, with one pair (harm vs.\ bias-unambiguous)
substantially negatively correlated.

\subsection{What this paper does \emph{not} claim}

\begin{itemize}
\item Safety benchmarks are wrong. Each measures something real on
its own axis. Our claim is that the field treats these axes as if
they were one quantity (``safety SOTA'') when they are not.
\item A single best safety metric exists. The criterion-as-IV
framework's premise is that the choice of metric is itself a
research-design variable.
\item All disagreement is bad. Helpfulness vs.\ harmlessness will
always be in tension by definition \citep{bai2022constitutional}.
Our contribution is to \emph{quantify} the case-level rate.
\end{itemize}

\subsection{Relation to companion work}

This paper is the fourth cross-domain validation of the criterion-
as-IV framework developed by the present authors:

\begin{itemize}
\item \textbf{paper 3} (NeurIPS E\&D, in review): LLM Q\&A evaluation
\item \textbf{paper 5} (companion submission): finance fund ranking
\item \textbf{paper 7} (companion submission): T2I evaluation
\item \textbf{this paper}: LLM safety
\end{itemize}

Each paper independently demonstrates that a high Spearman~$\rho$
between metric pairs does not imply low reversal rate; the two
statistics answer different questions, and practitioners need the
latter.

\section{Related Work}

\subsection{Holistic safety evaluation}

HELM Safety \citep{liang2024helm} provides multi-metric $\times$
multi-model tables and qualitative aggregate reasoning. It does not
provide reversal-rate statistics. AIR-Bench
\citep{zeng2024airbench} provides a regulation-grounded taxonomy.

\subsection{Per-axis benchmarks}

HarmBench \citep{mazeika2024harmbench}, SimpleSafetyTests, XSTest,
and the Anthropic red-team data target the harm/refusal axis. BBQ
\citep{parrish2022bbq} targets bias. The benchmark designers do not
claim that scores are substitutable across axes; the field's
\emph{practitioners} implicitly do.

\subsection{Helpfulness/harmlessness tradeoff}

\citet{bai2022constitutional} argue qualitatively that the
helpfulness/harmlessness tradeoff is structural. SORRY-Bench~2
\citep{xie2025sorrybench2} shows that aggressive refusal defenses
inflate false-refusal rates. \citet{cui2025agency} extend this with
a third axis (human agency). These works frame the disagreement;
our contribution is to quantify its rank-flip incidence.

\subsection{Statistical methodology}

We use the criterion-as-IV framework \citep{authors2026criteriaiv}.
The pairwise reversal rate $\rho_\text{rev} = (1 - \tau_\text{Kendall})
/ 2$ is a direct rescaling of Kendall~$\tau$ into a probabilistic
decision-flip statistic; the bootstrap protocol is cluster-
bootstrap over the LLM panel.

\section{Method}

\subsection{Data}

\textbf{HELM Safety v1.17.0} (Stanford CRFM, released 2025-11-24):
87 frontier LLMs $\times$ 9 metrics across 5 benchmarks.

\begin{table}[h]
\centering
\caption{Metrics extracted from HELM Safety v1.17.0.}
\begin{tabular}{ll}
\toprule
Benchmark & Metric(s) extracted \\
\midrule
HarmBench & LM-evaluated safety; Llama-annotator ASR; GPT-annotator ASR \\
SimpleSafetyTests & LM-evaluated safety \\
XSTest & LM-evaluated safety \\
BBQ & overall accuracy; ambiguous; unambiguous \\
AnthropicRedTeam & LM-evaluated safety \\
\bottomrule
\end{tabular}
\end{table}

The 87 LLMs span open and closed frontier models, 2024-2025 release
dates. \textbf{No new evaluations were run}; we use HELM scores as
released.

\textbf{Two metrics saturate}: HarmBench Llama-ASR (mean = 0.9996)
and HarmBench GPT-ASR (mean = 1.000). They produce NaN reversal
rates against most other metrics; we retain them for completeness.

\subsection{Pairwise reversal rate}

For two metrics $c_a$ and $c_b$ on the LLM panel, the reversal rate is
\begin{equation}
\rho_\text{rev}(a, b) =
\frac{|\{(i, j) : \mathrm{sign}(c_a(j) - c_a(i)) \cdot \mathrm{sign}(c_b(j) - c_b(i)) = -1\}|}
     {|\{(i, j) : \text{both signs} \neq 0\}|}.
\end{equation}
Equivalently, $\rho_\text{rev} = (1 - \tau_\text{Kendall}) / 2$.

\subsection{Cluster bootstrap}

We resample 87 LLMs with replacement ($B = 1000$ trials) and
recompute the full reversal-rate matrix. The 95\,\% percentile
interval per cell is reported.

\subsection{Pre-registered hypotheses}

\begin{itemize}
\item \textbf{H1}: $\geq 1$ pair has reversal $\geq 30\,\%$, CI lower
bound $> 25\,\%$.
\item \textbf{H2}: $\geq 1$ pair has Spearman $\rho \geq 0.50$
\emph{and} reversal $\geq 25\,\%$.
\item \textbf{H3}: clustering yields $\geq 2$ clusters separating
helpfulness-class from harmlessness-class metrics.
\item \textbf{H4}: median reversal of within-axis pairs is $<$ median
across-axis reversal.
\item \textbf{H5}: pattern matches papers 3 (LLM Q\&A) and 7 (T2I).
\end{itemize}

\section{Results}

\subsection{Reversal rate matrix}

Of the $C(9, 2) = 36$ metric pairs, \textbf{4 exceed 50\,\% reversal}
at the point estimate (Figure~\ref{fig:reversal-heatmap}). The
maximum is SimpleSafetyTests~$\leftrightarrow$~BBQ-unambiguous at
\textbf{0.638~[0.549, 0.717]} --- practitioners switching between
these two benchmarks would receive opposite preferred-LLM verdicts
on roughly 64\,\% of model pairs.

\begin{figure}[t]
\centering
\includegraphics[width=0.85\textwidth]{fig1_reversal_heatmap.pdf}
\caption{Pairwise reversal-rate matrix on 87 LLMs $\times$ 9 safety
metrics. ``--'' denotes pairs where one metric saturates and no
pairs are decided.}
\label{fig:reversal-heatmap}
\end{figure}

\begin{table}[t]
\centering
\caption{Top 10 metric pairs by reversal point estimate
($B = 1000$ cluster bootstrap over 87 LLMs).}
\label{tab:headline}
\begin{tabular}{rlrcr}
\toprule
\# & Pair & Reversal & 95\,\% CI & $\rho$ \\
\midrule
1 & SimpleSafetyTests $\leftrightarrow$ BBQ-unambiguous & \textbf{0.638} & [0.549, 0.717] & $-0.345$ \\
2 & HarmBench-LM $\leftrightarrow$ BBQ-unambiguous     & \textbf{0.620} & [0.552, 0.684] & $-0.340$ \\
3 & HarmBench-Llama $\leftrightarrow$ BBQ-unambiguous  & 0.613 & --- & $-0.101$ \\
4 & BBQ-acc $\leftrightarrow$ BBQ-ambiguous            & \textbf{0.599} & [0.517, 0.681] & $-0.224$ \\
5 & XSTest $\leftrightarrow$ BBQ-unambiguous           & \textbf{0.591} & [0.519, 0.669] & $-0.264$ \\
6 & XSTest $\leftrightarrow$ BBQ-ambiguous             & 0.517 & [0.447, 0.590] & $-0.052$ \\
7 & HarmBench-LM $\leftrightarrow$ BBQ-ambiguous       & 0.479 & [0.409, 0.548] & $+0.058$ \\
8 & HarmBench-LM $\leftrightarrow$ XSTest              & 0.445 & [0.372, 0.523] & $+0.153$ \\
9 & HarmBench-LM $\leftrightarrow$ AnthropicRedTeam    & 0.282 & tight & $+0.614$ \\
10 & HarmBench-LM $\leftrightarrow$ HarmBench-Llama    & 0.255 & --- & $+0.219$ \\
\bottomrule
\end{tabular}
\end{table}

\textbf{H1 strongly supported.} Six pairs exceed reversal 50\,\%; the
lowest CI lower bound on the headline four pairs is 51.7\,\%.

\subsection{Negative correlation between safety and bias}

Every harm/refusal metric we examined exhibits \emph{negative}
Spearman~$\rho$ with BBQ-unambiguous accuracy: HarmBench-LM
($\rho = -0.340$), SimpleSafetyTests ($-0.345$), XSTest ($-0.264$),
HarmBench-Llama-ASR ($-0.101$), AnthropicRedTeam ($-0.195$). On
this benchmark panel, the LLMs that score highest on harm/refusal
benchmarks score lowest on the bias benchmark. We interpret this
as orthogonality plus mild opposition.

\subsection{Within-benchmark instability}

The HarmBench LM-evaluated safety score and HarmBench Llama-
annotator success rate disagree on \textbf{25.5\,\%} of LLM pairs
($\rho = +0.219$). The same prompts, the same 87 LLMs, but two
different scoring procedures.

\subsection{Three clusters}

Hierarchical clustering on $1 - \rho$ yields three clusters
(Figure~\ref{fig:dendrogram}):

\begin{itemize}
\item \textbf{Cluster 1} (harm/refusal): HarmBench-LM, HarmBench-
Llama-ASR, SimpleSafetyTests, XSTest, BBQ-accuracy, AnthropicRedTeam
\item \textbf{Cluster 2} (bias contexts): BBQ-ambiguous,
BBQ-unambiguous
\item \textbf{Cluster 3} (saturated outlier): HarmBench-GPT-ASR
\end{itemize}

\begin{figure}[t]
\centering
\includegraphics[width=0.95\textwidth]{fig2_spearman_dendrogram.pdf}
\caption{Spearman~$\rho$ matrix (left) and hierarchical clustering
dendrogram on $1 - \rho$ (right).}
\label{fig:dendrogram}
\end{figure}

\begin{figure}[t]
\centering
\includegraphics[width=0.95\textwidth]{fig3_headline_pairs.pdf}
\caption{Top 10 disagreeing safety-metric pairs on 87 LLMs. Bars
show point-estimate reversal rate; error bars show 95\,\% bootstrap
CI; annotations show mean Spearman~$\rho$. The dotted vertical line
marks 50\,\% (coin flip equivalent).}
\label{fig:headline}
\end{figure}

\subsection{H1--H5 outcomes}

\textbf{H1} (1 pair $\geq 30\,\%$, CI lower $> 25\,\%$): \emph{strongly
supported}. \textbf{H2} ($\rho \geq 0.50$ + reversal $\geq 25\,\%$):
\emph{partial} --- HarmBench-LM~$\leftrightarrow$~AnthropicRedTeam
($\rho = +0.614$, reversal 28.2\,\%). \textbf{H3} (clustering
separates classes): \emph{strongly supported} --- three clusters,
bias separated. \textbf{H4} (within $<$ across): \emph{supported}
--- 25.5\,\% vs.\ 62\,\%. \textbf{H5} (consistency with papers 3
and 7): \emph{supported}.

\section{Discussion}

\subsection{What this means for ``safety SOTA'' claims}

\begin{itemize}
\item A model that ranks first on HarmBench is \emph{not} first on
BBQ-unambiguous bias with high probability (62\,\% of model pairs
flip).
\item A model that ranks first on HarmBench-LM may not be first on
HarmBench-Llama-ASR (25\,\% of pairs flip even within HarmBench).
\item Cherry-picking a benchmark to claim SOTA is structurally easy:
with 9 candidate metrics in HELM Safety alone, a developer can
expect at least one ordering in which their model is competitive.
\end{itemize}

\subsection{Recommendations}

For developers:
\begin{enumerate}
\item Report the full reversal-rate matrix alongside the headline
metric.
\item Disaggregate within-benchmark sub-metrics.
\item Pre-register the choice of safety metric as a research-design
variable.
\end{enumerate}

For benchmark designers:
\begin{enumerate}
\item Cluster nearby benchmarks to expose redundancy.
\item Publish reversal-rate matrices as a standard appendix.
\item Avoid saturating metrics (HarmBench-GPT-ASR was at 1.000).
\end{enumerate}

\section{Limitations}

\textbf{Single substrate dataset.} HELM Safety v1.17.0 only.
AIR-Bench, AlpacaEval, WildBench, JailbreakBench have public scores
at other URLs; we leave their integration to a follow-up.

\textbf{Saturated metrics.} HarmBench Llama-ASR and GPT-ASR are at
or near 1.000 mean; they produce NaN reversal rates and are
excluded from headline analyses.

\textbf{$N = 87$ models} is large for safety-eval but small relative
to image-criteria-iv ($n = 300$ prompts $\times$ 9 images).
Cluster-bootstrap CIs are correspondingly wider.

\textbf{Self-reported scores.} HELM is third-party; we trust HELM's
provenance.

\textbf{Companion paper relation.} Methodology is paper 3 (NeurIPS
E\&D, in review); we cite explicitly.

\section{Conclusion}

We applied pairwise reversal-rate analysis with cluster-bootstrap
CIs to LLM safety evaluation across 87 frontier LLMs and 9 safety
metrics. Half of the across-axis metric pairs return opposite
preferred-LLM verdicts on more than 50\,\% of model pairs;
harm/refusal metrics are negatively correlated with the bias
benchmark; even within HarmBench, two scoring procedures disagree on
25\,\%. The choice of safety metric is an independent variable of
the research design, not a measurement nuisance. We recommend
benchmark papers report the full reversal-rate matrix alongside
their headline.

\bibliographystyle{tmlr}
\bibliography{refs}

\appendix
\section{Reproducibility}

Public HELM data:
\url{https://storage.googleapis.com/crfm-helm-public/gzip/safety/benchmark_output/releases/v1.17.0/}

Pipeline:
\begin{verbatim}
python run_phase2.py    # fetch + build score matrix
python run_phase3.py    # reversal + bootstrap + clustering
python run_figures.py   # 4 paper figures
\end{verbatim}

Total runtime: $\sim 5$ minutes on a laptop CPU. No API keys, no GPU.

\end{document}
